% Artigo científico ABNT (NBR 6022) — abnTeX2
% Compile: pdflatex main && bibtex main && pdflatex main && pdflatex main

\documentclass[
  article,
  12pt,
  oneside,
  a4paper,
  english,
  brazil
]{abntex2}

\usepackage[T1]{fontenc}
\usepackage[utf8]{inputenc}

% Fonte: Times New Roman (padrão). Para Arial, comente TIMES e descomente ARIAL.
\usepackage{newtxtext,newtxmath}          % TIMES
% \usepackage[scaled=0.95]{helvet}        % ARIAL
% \renewcommand{\familydefault}{\sfdefault}

\usepackage{indentfirst}
\usepackage{graphicx}
\usepackage{microtype}
\usepackage[alf]{abntex2cite}

% Títulos na mesma família do corpo (Arial ou Times), não em sans-serif extra.
\renewcommand{\ABNTEXchapterfont}{\normalfont}

% --- Identificação (edite aqui) ---
\newcommand{\universidade}{Nome da Universidade}
\newcommand{\disciplina}{Nome da disciplina}
\newcommand{\professor}{Nome do professor}
\newcommand{\estudante}{Nome completo do estudante}
\newcommand{\matricula}{000000000}

\titulo{Título do artigo}
\tituloestrangeiro{Article title in English}
\autor{\estudante}
\instituicao{\universidade}
\local{Brasil}
\data{2026}

\makeatletter
\hypersetup{
  pdftitle={\@title},
  pdfauthor={\@author},
  colorlinks=true,
  linkcolor=black,
  citecolor=black,
  urlcolor=black
}
\makeatother

% Corpo 12pt, entrelinha 1,5, recuo de 1,25 cm, texto justificado.
% Margens 3 cm (superior/esquerda) e 2 cm (inferior/direita): padrão abnTeX2.
% Citações longas, notas e paginação: 10pt (\ABNTEXfontereduzida).
% Citações longas e referências: espaçamento simples (citacao / abntex2cite).
\OnehalfSpacing
\setlength{\parindent}{1.25cm}
\setlength{\parskip}{0pt}

\renewcommand{\imprimircapa}{%
  \begin{capa}
    \begin{center}
      {\large\bfseries\MakeUppercase{\universidade}\par}
      \vspace{1.5cm}
      Disciplina: \disciplina\par
      Professor(a): \professor\par
      \vspace{2.5cm}
      {\ABNTEXchapterfont\bfseries\LARGE\imprimirtitulo\par}
      \vfill
      {\large\estudante\par}
      Matrícula: \matricula\par
      \vfill
      {\large\imprimirlocal\par}
      {\large\imprimirdata\par}
    \end{center}
  \end{capa}
}

\begin{document}
\selectlanguage{brazil}
\frenchspacing

\imprimircapa
\maketitle

\begin{resumoumacoluna}
Texto do resumo em português (até cerca de 250 palavras).

\vspace{\onelineskip}
\noindent\textbf{Palavras-chave}: palavra1. palavra2. palavra3.
\end{resumoumacoluna}

\renewcommand{\resumoname}{Abstract}
\begin{resumoumacoluna}
\begin{otherlanguage*}{english}
Abstract text in English.

\vspace{\onelineskip}
\noindent\textbf{Keywords}: keyword1. keyword2. keyword3.
\end{otherlanguage*}
\end{resumoumacoluna}

\textual

\section{Introdução}

Texto introdutório. Exemplo de citação: \cite{exemplo2020}.

\section{Desenvolvimento}

Desenvolva o conteúdo do artigo aqui. Parágrafos do corpo usam tamanho 12,
espaçamento 1,5 e recuo de 1,25\,cm na primeira linha.

\begin{citacao}
Exemplo de citação direta longa (mais de três linhas): tamanho 10,
espaçamento simples e recuo em relação à margem. Substitua este bloco
pelo trecho citado \cite{exemplo2020}.
\end{citacao}

Continuação do texto após a citação longa.

\section{Considerações finais}

Conclusão breve.

\postextual
\bibliography{refs}

\end{document}
